\documentclass[12pt]{article}
 \usepackage[margin=1in]{geometry} 
\usepackage{amsmath,amsthm,amssymb,amsfonts, enumitem, fancyhdr, color, comment, graphicx, environ}
\usepackage{booktabs}
\usepackage[colorinlistoftodos]{todonotes}
\usepackage{algorithm}
\usepackage{algpseudocode}
\usepackage{tikz}
\usepackage{subcaption}
\usepackage{float}
\usetikzlibrary{arrows.meta, automata, positioning}

\usepackage{listings}
\usepackage{xcolor}
\lstset{
  basicstyle=\ttfamily\small,
  backgroundcolor=\color{gray!10},
  frame=single,
  breaklines=true,
  keywordstyle=\color{blue},
  commentstyle=\color{gray},
  stringstyle=\color{red},
  showstringspaces=false,
  numbers=left,            
  numberstyle=\tiny\color{gray}, 
  numbersep=10pt 
}

%%%%%%%%%%%%%%%%%%%%%%%%%%%%%%%%%%%%%%%%%%%%%%%%%%%%%%%%%%%%%%%%%%%%%%%%
\pagestyle{fancy}
\setlength{\headheight}{65pt}
\newenvironment{problem}[2][Problem]{\begin{trivlist}
\item[\hskip \labelsep {\bfseries #1}\hskip \labelsep {\bfseries #2.}]}{\end{trivlist}}
\newenvironment{sol}
    {\emph{Solution:}
    }
    {
    \qed
    }
\specialcomment{com}{ \color{blue} \textbf{Comment:} }{\color{black}} %for instructor comments while grading
\NewEnviron{probscore}{\marginpar{ \color{blue} \tiny Problem Score: \BODY \color{black} }}

\newcounter{subproblem}
% \renewcommand{\thesubproblem}{\alph{subproblem}} % letters 
\renewcommand{\thesubproblem}{\arabic{subproblem}} % numbers
\newenvironment{subprob}[1][]{
  \refstepcounter{subproblem}
  \begin{trivlist}
  \item[\hskip \labelsep {\bfseries (\thesubproblem)}]
}{
  \end{trivlist}
}
\newenvironment{subsol}
    {\emph{Solution:}
    }
    {
    \qed
    }

\newtheorem{theorem}{Theorem}
\newtheorem{lemma}{Lemma}
\newtheorem{proposition}{Proposition}
\newtheorem{corollary}{Corollary}

% numbered within sections (optional)
% \newtheorem{theorem}{Theorem}[section]
% \newtheorem{lemma}[theorem]{Lemma} % shares theorem counter

%%%%%%%%%%%%%%%%%%%%%%%%%%%%%%%%%%%%%%%%%%%%%%%%%%%%%%%%%%%%%%%%%%%%%%%%%%%%%%%%%
\setlength {\marginparwidth }{2cm}
\newcommand{\nextline}{\\[0.2em]}


\newtheorem*{remark}{Remark}

\newcommand{\R}{\mathbb{R}}
\newcommand{\N}{\mathbb{N}}
\newcommand{\Q}{\mathbb{Q}}
\newcommand{\Z}{\mathbf{Z}}
\newcommand{\st}{\text{ s.t }}
\newcommand{\tand}{\text{ and }}
\newcommand{\tor}{\text{ or }}
\newcommand{\bigO}[1]{\mathcal{O}\left(#1\right)}
\newcommand{\fP}{\mathbb{P}}
\newcommand{\E}{\mathbb{E}}


\usepackage{listings}
\usepackage{xcolor}
\lstset{
  basicstyle=\ttfamily\small,
  backgroundcolor=\color{gray!10},
  frame=single,
  breaklines=true,
  keywordstyle=\color{blue},
  commentstyle=\color{gray},
  stringstyle=\color{red},
  showstringspaces=false,
  numbers=left,            
  numberstyle=\tiny\color{gray}, 
  numbersep=10pt 
}


% titel en auteur van het document
\title{Theory of Algorithms Problem Set \#3}
\author{Nicole Tian}
% \date{} % de datum van het document 

% einde definities; start van de tekst
\begin{document}

\maketitle                  % maakt een kopje met titel, auteur en datum
\thispagestyle{empty}       % zet paginanummering uit voor deze pagina

% kopje met bijbehorende titel; varianten: \subsection, \subsubsection 
\section*{Problem 1: Recurrence Relations}
\begin{enumerate}
\item[(a)] $T(n) = 2T(\frac{n}{4}) + 2n = 4n - \sqrt{n}$
\begin{align*}
  T(n) &= 2T(\frac{n}{4}) + 2n \\
  &= 4T(\frac{n}{16}) + n + 2n \\
  &= 8T(\frac{n}{64}) + \frac{n}{4} + n + 2n \\
  &\vdots \\
  &= 2^kT(\frac{n}{4^k}) + \sum_{i=1}^{k}2^i(\frac{n}{4^{i-1}}) \\
  &= 2^kT(\frac{n}{4^k}) + 2n\sum_{i=0}^{k-1}(\frac{1}{2})^i \\
  &= 2^kT(\frac{n}{4^k}) - 4n(\frac{1}{2}^{k} - 1) \\
\end{align*}
We set $k = log_4(n)$
\begin{align*}
  &= \sqrt{n} * T(1) - 4n(\frac{1}{\sqrt{n}} - 1) \\
  &= 3\sqrt{n} - 4n(\frac{1}{\sqrt{n}} - 1) \\
  &= 3\sqrt{n} - (\frac{4\sqrt{n}\sqrt{n}}{\sqrt{n}} - 4n) \\
  &= 3\sqrt{n} - (4 \sqrt{n} - 4n)\\
  &= 4n - \sqrt{n}
\end{align*}

\item[(b)] $T(n) = 4T(\frac{n}{3}) + 5n^2 = 9n^2 - 8n^{log_3(4)}$
\begin{align*}
  T(n) &= 4T(\frac{n}{3}) + 5n^2 \\
  &= 4^kT(\frac{n}{3^k}) + 5n^2\sum_{i=0}^{k-1}(\frac{4}{9})^i \\
  &= 4^kT(\frac{n}{3^k}) + 5n^2\frac{9}{5}(1-(\frac{4}{9})^k)
\end{align*}
We set $k = log_3(n)$
\begin{align*}
  &= n^{log_3(4)} T(1) + 9n^2(1-\frac{n^{log_3(4)}}{n^2}) \\
  &= n^{log_3(4)} + 9n^2 - 9n^{log_3(4)}\\
  &= 9n^2 - 8n^{log_3(4)}\\
\end{align*}


\item[(c)] $T(n) = T(\frac{n}{2}) + 2n - 1 = 4n - log_2(n) + 3$
\begin{align*}
  T(n) &= T(\frac{n}{2}) + 2n - 1 \\
  &= T(\frac{n}{2^k}) - k + 2n\sum_{i=0}^{k+1}(\frac{1}{2})^i \\
  &= T(\frac{n}{2^k}) - k + 4n(1-(\frac{1}{2})^k) \\
\end{align*}
We set $k = log_2(n)$
\begin{align*}
  &= T(1) - log_2(n) + 4n(1-\frac{1}{n}) \\
  &= 1 - log_2(n) + 4n - 4 \\
  &= 4n - log_2(n) + 3 \\
\end{align*}
\item[(d)] 5 points extra credit: $T(n) = 12T(\frac{n}{2}) - 32T(\frac{n}{4}) + n^2 + 3n$, $T(1) = 1$, $T(2) = 2$.
\end{enumerate}

\section*{Problem 2}
Consider the following situation. There is a problem you are trying to design an algorithm for. For this problem, you have a way to evaluate the performance of any algorithm at each time $t$. In particular, for any algorithm $A$, you have a function $v_A: \mathbb{N} \cup \{0\} \to \mathbb{R}$ which gives the performance $v_A(t)$ of the algorithm at time $t$ (where we always set $v_A(0) = 0$ as there is no progress at time $0$). 

You have designed a greedy algorithm $A$ and you think you have an argument for why it should always stay ahead. However, your partner comes up with another algorithm $B$ and you find that at the end, algorithm $B$ did better than algorithm $A$, i.e., $v_{B}(t_{end}) > v_{A}(t_{end})$. 

You would like to figure out where your reasoning failed. To do this, you would like to find a time $t \in \{0,1,\ldots,t_{end}-1\}$ such that $v_{A}(t) \geq v_{B}(t)$ but $v_{A}(t+1) < v_{B}(t+1)$, i.e., your greedy algorithm was ahead (or at least equal) at time $t$ but fell behind at time $t+1$.

For example, if $t_{end} = 8$ and the performances of the algorithms were as follows
\begin{center}
\begin{tabular}{|c|c|c|c|c|c|c|c|c|c|c|c|}
 \hline
 Time $t$ & $0$ & $1$ & $2$ & $3$ & $4$ & $5$ & $6$ & $7$ & $8$\\
 \hline
 $v_A(t)$ & $0$ & $7$ & $9$ & $10$ & $14$ & $18$ & $24$ & $28$ & $30$\\
 \hline
 $v_B(t)$ & $0$ & $4$ & $7$ & $8$ & $16$ & $17$ & $22$ & $30$ & $32$\\
 \hline
\end{tabular}
\end{center}
then either $t = 3$ or $t = 6$ would be a correct answer.
\begin{enumerate}
\item[(a)] 10 points: Give a fast algorithm for finding such a time $t$.
\item[(b)] 5 points: Explain why your algorithm is correct. 
\item[(c)] 5 points: What is the runtime of your algorithm? Explain why your algorithm has this runtime.
\end{enumerate}



\section*{Problem 3: Target Sum (25 points)}
You have $n$ integers $x_1,\ldots,x_n$ and you want to either find four of these integers which sum to a target integer $S$ or show that this is impossible. In other words, you either want to find four indices $i < j < k < l \in [n]$ such that $x_i + x_j + x_k + x_l = S$ or show that no such indices exist.
\begin{enumerate}
\item[(a)] Consider the following algorithm:
\begin{enumerate}
  \item[1)] Create 3-tuples of the form $(V, i, j)$ with $V$ representing the combined value of integers $x_i$ and $x_j$ s.t. $i < j$. 
  \item[2)] Sort by $V, i, j$.
  \item[3)] perform two-pointer search to look for two values $V$ that add up to $S$.
  \item[4)] If we find two such values, we iterate through all the right 3-tuples that have value $V_R$ and create two maps from $i$ to $j$ and vice versa. For each 3-tuple $(V, i, j)$, we add $j$ to i-map[i] and $i$ to j-map[j]. Note down the total number of 3-tuples. 
  \item[5)] Move left-pointer along the 3-tuples of value $V_L$. We take each 3-tuple, $(V_L, i_l, j_l)$. 
  \begin{enumerate}
  \item The number of tuples with a value of $V_R$ which contains either $j_l$ and $i_l$ can be represented as follows with \\
  
  $\sigma=
    \begin{cases}
      1, & V_R = V_L\\
      0, & V_R \neq V_L
    \end{cases}$
  
    \begin{equation}
      |\text{Tuples}_{i_l, j_l}| = |\text{i-map}[i_l]| + |\text{i-map}[j_l]| + |\text{j-map}[i_l]| + |\text{j-map}[j_l]| - \sigma
    \end{equation}
    
    If $|\text{Tuples}_{i_l, j_l}| < |\text{Tuples}|$, then we have a solution and we can iterate through the values to look for it. 
  \item Otherwise, we move onto the next tuple on the left-side until there are no more 3-tuples of value $V_L$.
  \item If there exists no more 3-tuples of value $V_L$, we continue with the two pointer search algorithm with different values of $V_L$ and $V_R$ and repeat steps 2-5. \\
  
  Note that if $V_L + V_R = S$, then we cannot have another $V_L' + V_R = S \st V_l' \neq V_L$, so we move onto completely new values of $V_L$ and $V_R$, and therefore, under this algorithm, we never rebuild the $i$-$j$-mappings for any value $V$ more than once nor do we traverse any tuple twice. 
  \end{enumerate}
  \item[6)] If the two-pointers meet each other without finding a solution, we are done and no solution exists. 
\end{enumerate}

\item[(b)] We will prove the above algorithm is correct by proving that if there exists a solution, the algorithm will find one, and if the algorithm finds no solution, there does not exist a solution.
\begin{lemma}
  Denote $R$ to be the set of all 3-tuples of value $V_R$ and $L$ to be the set of all 3-tuples of value $V_L$. 
  
  The algorithm outputs a solution $\iff$ there exist $(V_L,i,j)\in L$ and $(V_R,k,l)\in R \st \{i,j\}\cap\{k,l\}=\varnothing$ 
  
  Note that, the algorithm is not guarneteed to return any specific pairs.
\end{lemma}
\begin{proof} $\nextline$
  Fix $V_L$ and $V_R \st V_L+V_R=S$ and consider the right block $R$. For any indices $i,j$, let $T(i,j)$ denote the number of tuples in $R$ that contain either $i$ or $j$. Then the number of tuples in $R$ that contain neither $i$ nor $j$ is $|R|-T(i,j)$.

  By construction of the maps in step (4), the algorithm computes
  \[
  T(i,j)=|\text{i-map}[i]|+|\text{j-map}[i]|+|\text{i-map}[j]|+|\text{j-map}[j]|-\sigma,
  \]
  where $\sigma=1$ if $V_L=V_R$ (equivalently, $(i,j)\in R$) and $\sigma=0$ otherwise.
  Thus the algorithm tests whether $T(i,j)<|R|$, which is equivalent to
  $|R|-T(i,j)>0$.

  ($\Rightarrow$) If the algorithm outputs a solution, then for some $(V_L,i,j)\in L$ it found $T(i,j)<|R|$, hence $|R|-T(i,j)>0$. Therefore there exists a tuple $(V_R,k,l)\in R$ that contains neither $i$ nor $j$, i.e.\ $\{i,j\}\cap\{k,l\}=\varnothing$, so disjoint pairs exist.

  ($\Leftarrow$) If there exist disjoint pairs $(V_L,i,j)\in L$ and $(V_R,k,l)\in R$ with $\{i,j\}\cap\{k,l\}=\varnothing$, then at least one tuple in $R$ contains neither $i$ nor $j$.
  Hence $|R|-T(i,j)>0$, equivalently $T(i,j)<|R|$, so the algorithm will pass the test in
  step (5a) and output a solution.
\end{proof}

\begin{theorem}
  There exists a solution $\iff$ the algorithm finds a solution
\end{theorem}
\begin{proof} $\nextline$
  ($\Rightarrow$)
    Assume $\exists x_i, x_j, x_k, x_l \st x_i + x_j + x_k + x_l = S \tand i < j < k < l$. $\nextline$
    WLOG, assume that $x_i < x_j < x_k < x_l$. Then $x_i + x_j < x_k + x_l$.
    By Step (1) and (2) in our algorithm, $\exists V_1= x_i + x_j \tand V_2 = x_k + x_l$. $\nextline$
    Our pointers start out on $V_L \leq V_1, \tand V_R \geq V_2$. By the algorithm, the $V$-values in sorted increasing order and therefore, the two pointer algorithm will find a solution to $V_L + V_R = S$ if it exists. $\nextline$
    A solution exist by assertion and therefore we will find one. Note this solution may not be $V_1, V_2$ and so we will proceed with some arbitrary values $V_L, V_R$. $\nextline$
    Within $V_L, V_R$, there are two possibilities: $\exists$ two disjoint pairs or $\nexists$ exists two disjoint pairs. If the disjoint pairs exist, by Lemma 1, the algorithm will find it and the proof ends trivially. If the disjoint pairs do not exist then we know that $V_L \leq V_1, \tand V_R \geq V_2$ and we continue the algorithm, discarding the current $V_L, V_R$. This is because two-pointer starts with the most extreme pair of high and low values so $V_1 \tand V_2$ must be more moderate than $V_L \tand V_R$ if $V_L \tand V_R$ were found first. $\nextline$
    Therefore, we will either land on some $V_L, V_R$ which contains two disjoint pairs and the proof ends, or we will land on $V_1, V_2$ and because it is guarenteed to contain two disjoint pairs in $V_1 \tand V_2$, by Lemma 1, the algorithm will output a solutions.

  ($\Leftarrow$) Assume the algorithm finds a solution. By Step (3) of our algorithm, it means that ours pointers found some $V_L, V_R \st V_L + V_R = S$. By Lemma 1, there exist $(V_L,i,j)\in L$ and $(V_R,k,l)\in R \st \{i,j\}\cap\{k,l\}=\varnothing$. Therefore $x_i + x_j + x_k + x_l = V_L + V_R = S$
\end{proof}

\item[(c)] Consider the time complexity of each step in the algorithm:
\begin{enumerate}
  \item[1)] Creating pair sums of the form $(V, i, j)$ creates $\binom{n}{2} \text{ pairs } \Rightarrow O(n^2)$ time
  \item[2)] Sorting the 3-tuples is a mergesort with $n^2$ objects $\Rightarrow O(n^2log(n))$
  \item[3)] Traversing the 3-tuples $\Rightarrow O(n^2)$ time
  \item[4)] Creating the mappings for every value of $V$ takes $O(n^2)$ overall because we add each 3-tuple to some mapping once
  \item[5)] Iterating within the mapping if there exists a solution happens at most once and in the worst case takes $O(n^2)$ comparisons
\end{enumerate}
$\Rightarrow O(n^2) +O(n^2log(n)) + O(n^2) + O(n^2) + O(n^2) = O(n^2log(n))$ 

\end{enumerate}




\section*{Problem 4: Rephrasing of Problem 5 in Chapter 5 of Kleinberg-Tardos (25 points)}
Let's say that we are given $n$ lines $\{y = {a_i}x + b_i: i \in [n]\}$ and we want to find all of the lines which can be seen from above without passing through another line. In other words, we want to find all $i \in [n]$ such that $\exists x \in \mathbb{R} \left(\forall j \in [n] \setminus \{i\} ({a_i}x + b_i \geq {a_j}x + b_j)\right)$.

\begin{enumerate}
\item[(a)] Denote all lines, $\{y = {a_i}x + b_i: i \in [n]\}$, as tuples of the form $(a_i, b_i)$. We sort the tuples by $a_i$. We iteratively go through the tuples, start with the second tuple, denoted as $(a_c, b_c)$:
\begin{enumerate}
  \item[(1)] Take the tuple with an index one less than the current, denote as $(a_p, b_p)$, if it exists, otherwise continue onto the next tuple
  \item Compute intersection of the two tuples as lines. Denote this point as $p_c$
  \item[(2)] Take the second most recently computed point, $p_p$, if it exists. Otherwise, continue onto the next tuple 
  \item[(3)]Compare $p_c$ and $p_p$. If $p_c > p_p$, we continue onto the next tuple. Otherwise, we remove $(a_p, b_p)$ and $p_p$ and repeat step 1-4.
  \item[(4)] The final answer is the number of tuples left, or the number of points + 1.
\end{enumerate}
\item[(b)] 10 points: Explain why your algorithm is correct.
\item[(c)] 5 points: What is the runtime of your algorithm? Explain why your algorithm has this runtime.
\end{enumerate}
\newpage
\section*{Hints}
3a. You should consider sums of pairs of the integers and how they can be combined. When trying to combine these sums, be careful that you do not use the same integer twice.\\
\\
4a. You can divide up the lines based on their slope. If you solve the problem for the subset of lines with smaller slopes and the subset of lines with larger slopes, can you combine your solutions to obtain a solution for the entire problem?

\end{document}
